\section{Integration: What Makes a Self}

A higher vibe is a coherent crowd. But what separates a genuine crowd-that-is-a-self from a mere heap, a pile of unrelated vibes that happens to be nearby? A rock is a crowd of vibes too, and so is a cloud of gas, yet neither is a self. The answer is a single measurable quantity, integration, and this chapter explains it.

\subsection{The test: can you cut it without loss?}

Integration measures how much a crowd is bound into one, how hard it is to cut into independent pieces without destroying something. A genuine self resists cutting. Slice it in half and you have not made two smaller selves. You have wrecked it, because its parts were deeply working together, each depending on the rest. A mere heap, by contrast, cuts cleanly. Split a pile of sand and you just have two piles of sand. Nothing was lost, because the grains were never really cooperating.

So the question that defines a self is: how much does this whole exceed the sum of its parts? How much would be lost if you cut it apart? A lot, for a self. Almost nothing, for a heap. That ``how much would be lost when cut'' is integration.

\subsection{Measuring it on the crystal}

On the crystal this becomes a precise number. Take a region and ask how strongly its vibes are woven together through their notes, how much information they share that no subset holds alone. A tightly woven cluster scores high. A loose scatter scores near zero. The simulation computes this directly: a cohesive cell of the body scores high integration, a random same-size bag of vibes scores essentially zero, and a genuine self sits at a local peak, meaning if you add outsiders or remove core members, the integration drops. The self is the locally most-woven region, the natural unit, not an arbitrary boundary you drew.

This is satisfying because it means selfhood is not in the eye of the beholder. There is a fact of the matter about where one self ends and another begins: the boundaries are where integration peaks. Nature carves itself at these joints.

\subsection{One dial doing three jobs}

The beautiful thing is that this one quantity, integration, turns out to govern several things at once.

It decides what counts as a being: high integration means a genuine self, low means a heap. It decides whether the higher-level law holds: recall from the last chapter that the clean higher rule appears exactly when you coarsen along integrated wholes. And, as the next chapters show, it decides whether a disturbance gets absorbed and healed or instead takes on a life of its own. One dial, three jobs, which is a sign the theory found a real quantity rather than three unrelated ones.

\subsection{A bridge to consciousness research}

This idea is not invented from scratch. It echoes a serious line of consciousness research (integrated information theory) that proposes consciousness is tied to exactly this property, how much a system is unified and irreducible. The crystal gives that idea a concrete home: integration is computable on the web, it picks out the selves, and the most integrated structures are the candidates for the richest experience. The theory does not claim to have solved consciousness, but it places the leading structural measure of it at the center, where it does real work.

A caution, kept for later but flagged now: measuring how integrated a self is, is not the same as settling whether, or how, it feels. Integration tells you where the unified wholes are and how unified they are. Whether that structural unity is the same as a single felt experience is the deep open question we will name honestly. Integration is the structure of selfhood. The feeling of selfhood is the territory that structure maps.

\textbf{Takeaway:} integration measures how much a crowd is bound into one, how much would be lost if you cut it apart. It is high for a genuine self, near zero for a heap, and peaks exactly at a self's natural boundary. This one measurable dial decides what is a being, where the higher law holds, and how disturbances are handled, and it connects the theory to serious consciousness research.
